% エントリ 18(2026-11-11): χ_SR を基準の運動論の表にし、精度の約束を決めて全範囲で照合した
\section{[E18, 2026-11-11] χ$_\mathrm{SR}$ を基準の運動論の表にし、精度の約束を決めて入力の全範囲で照合した}\label{sec:e18}

共同研究者に渡す χ$_\mathrm{SR}$ について、「人に渡すなら、正しいと確信が持てるものを渡す」(ユーザー)。E17 までの関数は、半解析の各部分を ODE と点で比べただけで、$f>10^{17}\GeV$ の 322 の段階のように 0.19 ずれる所が残っていた。そこで (1) 基準の運動論を一つ決め、(2) 精度の約束をユーザーが決め(2026-11-10、「0.03、細かい表」)、(3) 値をその運動論から直接作った表で引くように作り直し、(4) 表を作るのに使っていない点で、入力の全範囲で照合した。計算は \texttt{calc/fp3.py}(\texttt{evolve3m})、\texttt{calc/reftab.py}、\texttt{calc/e18\_reftab\_build.py}、\texttt{calc/e18\_validate2.py}、単体版は \texttt{calc/export\_chisr.py}。

\subsection{約束}

\textbf{印 0・1・2 の値は、基準の運動論と $|\chi_\mathrm{SR}(T)-\chi_\mathrm{ref}(T')|\le0.03$ を満たす $T'$ が $|\ln(T'/T)|\le0.02$ の中にある。}
速い段差(雲が飽和してスピンが急に落ちる所)の時刻は 2\% までずれてよい、という意味である。0.03 は GW のスピンの測定誤差($\sim0.1$)より十分小さい。約束を満たせない所は印をつける: 印 4(模型の外の準位が育ちうる、E15・E17)と、新しい印 5(下)。

\subsection{基準の運動論}

\texttt{evolve3m} は 211・322・544 の占有数($\ln n_i$)と BH の質量・角運動量を、1 個の量子から時間で解く。
\begin{itemize}
\item 超放射は相対論的な率の表(E4、E16)。BH は $dM=-(\omega_i/\mu)\alpha_\odot\,dn_i$、$dJ=-m_i\,dn_i$ を失う。
\item 自己相互作用は BGLS の Table V・VI の率。BH が吸収する過程 $i\times j\to k\times\mathrm{BH}$ では、BH がエネルギー $\omega_i+\omega_j-\omega_k$ を受け取る($m$ は保存され、角運動量は変わらない)。
\item 雲からの重力波(BGLS Table IV): 211 と 322 の消滅($10^{-2}\alpha^{14}$、$3\times10^{-8}\alpha^{18}$)と 322 $\to$ 211 の遷移($5\times10^{-6}\alpha^{10}$)。自由場では、322 の段階のあいだに残った 211 の雲が BH に吸い戻されて回転を戻すのを、この消滅が防ぐ。
\item 減衰する準位は占有数を $10^{-2}$ 個の量子で止める(解法が $\ln n\to-\infty$ に追われないため)。解法は LSODA($\mathrm{rtol}=10^{-8}$)で、止まったら最後の出力の時刻から Radau、BDF に切り替えて続ける。
\item 確かめ: $\infty$ への過程を切ると、角運動量は $2\times10^{-6}$、エネルギーは $7\times10^{-6}$ の精度で保たれる。
\end{itemize}
率はすべて $\mu$ に比例するので時間は $\tau=\mu[\mathrm{eV}]\,t[\mathrm{yr}]$ で入り、$M$ は 1 個の量子から数える $e$-fold の数 $L=\ln(N_0M^2)$ を通してだけ効く。

\subsection{表の形}

\begin{itemize}
\item \textbf{座標}: 臨界に対する位置 $w$($w=0$ が 322 の臨界、1 が 211 の臨界、2 が 0.995)。初期の位置 $s$ で、322 から育つ側($s<1$、$\sigma=\ln s$)と 211 から育つ側($\sigma=\ln(s-1)$)に分ける。スピンの進み具合 $b$ は、最初の段を $[0,1]$ に、臨界より下は $w$ そのもので測る。
\item \textbf{進み具合の時刻}: スピンは最大 0.03 戻ることがある(BH の質量の増加と 211 の吸い戻し)。そこで「その時刻より後の $b$ の最小」(最後に通る包絡線)を進み具合とし、$b=0$--4 の 0.01 刻み(401 準位)を最初に越えた $\log_{10}\tau$ を持つ。包絡線からの戻り $b-b_\mathrm{lp}$ は $\log_{10}\tau$ の 0.05 dex 刻みで持つ。最初に越えた時刻や行き過ぎの包絡線では、行き過ぎの深さが節ごとに違って準位の時刻がそろわなかった(試作で確かめた)。
\item \textbf{節}: $\alpha$ 0.008--0.5 の対数で 47 点、$\sigma$ 25 点、$\log_{10}f$ 8--20.5 の 0.25 刻みで 51 点、$M$ 1・300$\Msun$。239,700 節を 10 並列で 18,904 秒(Mac Studio)。235,163 節が最後まで解け、1,987 節は途中で止まった(止まった $\tau$ まで準位を持ち、その 1/10 より先を引く点は印 3)。2,550 節は 211 の臨界が 0.995 を越えるので作らない。
\item \textbf{引き方}: $(\log\alpha,\sigma,\log f)$ について準位ごとの $\log\tau$ を多重線形に補間し、$L$ について $\tau$ を線形に混ぜる。問いの $\tau$ を挟む準位を二分法で探す。$f>10^{20.5}\GeV$ は上端の値、$f<10^8\GeV$ と $\alpha<0.008$ は E17 までの半解析の値(どちらもほぼ何も起きない領域)。表は $\log\tau$ を uint16、戻りを float16 に丸めて持つ(単体版の表は 12 MB)。
\end{itemize}

\subsection{印 5: 段差の時刻が節のあいだで飛ぶ所}

322 の育ち方の機構が切り替わる所などで、段差の時刻が $f$ について 0.1 dex のうちに約 1.2 dex 飛ぶ(例: $\alpha=0.161$、$\log_{10}f\approx17.9$)。どの表でも節のあいだは補間できないので、印をつける。角の節それぞれの問いの時刻での値の幅が $\theta_v$ を越え、しかも角の節がいる準位で段差の時刻が両隣を結ぶ直線から $\theta_k$ [dex] 以上ずれる(曲がりの指標、4 準位ずつの最大)セルを印 5 とする。値は補間のまま返す。

閾値は、照合とは別の 2000 点(seed 1801)で決めた。$\theta_k\le0.07$ かつ $\theta_v\le0.10$ なら約束を越えるセルはない。$\theta_k=0.1$ にすると表の角($\alpha=0.498$、$f=2.3\times10^{20}\GeV$、544 の臨界まで下がる所)で 0.055、$\theta_v=0.12$ にすると GW241011 の層の 1 点($\alpha=0.11$、$f\approx10^{18}\GeV$、$1.3\times10^7\yr$)で 0.038 を越える。両方の端から余裕をとって $(\theta_v,\theta_k)=(0.06,0.05)$ とした。粗い表(節を一つおき)では、約束を守ると GW241011 の層の 13\% 以上に印 5 がついた。細かい表では 7\%。

\subsection{照合}

表を作るのにも閾値を決めるのにも使っていない 3000 点(seed 1802、4 つの層 × 750 点)× 40 の年齢($10^5$--$1.4\times10^{10}\yr$)= 12 万セル。
\begin{itemize}
\item 層: 一様($M$ 2--200$\Msun$、$\chi_0$ 0--0.995、$\mu$ $10^{-13}$--$10^{-11}\eV$、$f$ $10^6$--$10^{22}\GeV$)、臨界の近く(同じ範囲で $\chi_0$ を 211 か 322 の臨界から $10^{-4}$--0.05 に置く)、1 世代目($M$ 8--50$\Msun$、Beta(0.69, 2.28)、$\mu$ 0.6--1.0$\times10^{-12}\eV$、$f$ $10^{10}$--$10^{19}\GeV$)、GW241011($M$ 15--25$\Msun$、$\chi_0$ 0.6--0.95)。$\alpha\le0.5$。
\item 基準の運動論は年齢の 2400 点の格子で解き、年齢の窓の中の $\chi_\mathrm{ref}$ の幅と比べる。最後まで解けなかった点は、評価の回数 4 倍、$\mathrm{rtol}=10^{-7}$、$10^{-6}$ の順に解き直した。緩い rtol の曲線は、それまでの曲線と重なる年齢で $3\times10^{-3}$ 以内に合うときだけ先の年齢に使う($\mathrm{rtol}=10^{-6}$ は急な段差で $10^{-2}$ ずれた例がある)。最後まで解けなかったのは 1 点($\alpha=0.464$、$f=1.2\times10^{17}\GeV$、$2.7\times10^8\yr$ まで)で、照合できない 14 セルはすべて印 3・4・5 だった。
\end{itemize}

\begin{table}[H]
\centering\small
\begin{tabular}{lcc}
\toprule
 & seed 1801(閾値を決めた) & seed 1802(独立)\\
\midrule
印 0・1・2 のセル & 69,004 & 104,941\\
約束の誤差の最大 & 0.021 & 0.023\\
約束を越えたセル & 0 & 0\\
99\% 点 & 0.005 & 0.005\\
時刻のずれを許さない誤差の最大 & 0.022 & 0.050\\
照合できなかった印 0・1・2 & 0 & 0\\
印 0・1・2・3・4・5 の割合 [\%] & 27.7、37.6、20.9、6.6、4.7、2.4 & 27.5、40.1、19.9、6.6、3.5、2.5\\
印 5 の割合: 1 世代目、GW241011 [\%] & 0.7、7.0 & 0.7、6.9\\
印 4 の誤差(0.03 を越える割合、最大) & 0\%、0.013 & 0.4\%、0.17\\
印 5 の誤差(0.03 を越える割合、最大) & 3.0\%、0.27 & 3.0\%、0.26\\
\bottomrule
\end{tabular}
\caption{χ$_\mathrm{SR}$ と基準の運動論の照合。誤差は年齢の 2\% のずれを許した値。}
\label{tab:e18val}
\end{table}

seed 1802 の印 0・1・2 の約束の誤差の最大を範囲ごとに見ると、$f$ では $10^6$--$10^8\GeV$ で 0(半解析の領域)、$10^8$--$10^{16}\GeV$ で 0.008 以下、$10^{16}$--$10^{18}$ で 0.015、$10^{18}$--$10^{20}$ で 0.023、$10^{20}$--$10^{22}$ で 0.008(上端の値の領域)。$\alpha$ では 0.1 未満で 0.008 以下、0.1--0.3 で 0.015、0.3--0.4 で 0.023、0.4--0.5 で 0.005。

\begin{result}
\begin{itemize}
\item 独立な 12 万セルで、印 0・1・2 の約束の誤差は最大 0.023、越えたセルはない。どの範囲にも偏りはない。
\item 印 5 は全体の 2.5\%、1 世代目の層で 0.7\%、GW241011 の層で 6.9\%。GW241011 の層の印 5 の 8 割は $f\ge10^{16}\GeV$ にあり、窓の $10^{11}$--$10^{14}\GeV$ にあるのは 1 割。
\item 単体版はプロジェクト版と、無作為な 20 万個と照合の 12 万個の入力すべてで、値も印も完全に一致する(書き出しの確かめに加えた)。照合の結果は単体版にそのまま当てはまる。
\item E17 の限界 3($f>10^{17}\GeV$ の 322 の段階で最大 0.19)と、質量の減少・重力波の扱いの違い(領域の境目で 0.016 の飛び)はなくなった。
\end{itemize}
\end{result}

\begin{caveat}
\begin{itemize}
\item 約束は基準の運動論(三準位)への忠実さである。433 の降下(印 3・4 で示す)、自己相互作用の率の相対論的な補正(E7)は入っていない。
\item 固定点で $\alpha>0.30$ の点は、表は $\alpha\le0.5$ まであるが、固定点の判定の表がこの範囲なので印 3 のまま。
\item 表の分解能: 進み具合は段の幅の 1\% 刻みで、最初の刻みを越えるまでは初期のスピンを返す。表の節の上でも基準の運動論と 0.006 ずれる(書き出しの確かめの閾値は 0.01)。
\item 照合した範囲は $M$ 2--200$\Msun$、$\mu$ $10^{-13}$--$10^{-11}\eV$。範囲の外でも値は返す。
\item 速さは $10^6$ 個で 3.6 秒(1 世代目らしい入力)--9.2 秒($\chi_0$ が一様)と、E17 の 0.8 秒より遅い。$10^6$ 個を一度に渡すと約 3 GB を使う。
\item E10--E17 のスクリプトの \texttt{chisr.build()} も表を使うようになった(表は git に入れていないので、ないときは止まる)。過去のエントリの数値は半解析の版での値である。窓の数値(E12・E13)を表の版で計算し直すのは未着手。
\end{itemize}
\end{caveat}
